%! AP Statistics — Unit 8, Lesson 8.5 — Lesson Plan
\documentclass[10pt]{article}
\usepackage{apstats-article}
\usepackage{apstats-boxes}

\newcommand{\UnitNumberName}{Unit 8: Inference for Categorical Data: Chi-Square \quad}
\newcommand{\LessonNumberName}{Lesson 8.5: Setting Up a Chi-Square Test for Homogeneity or Independence}

\begin{document}
\begin{center}
    {\Large\bfseries \CourseName: \SchoolYear \\
    {\small\bfseries \UnitNumberName \LessonNumberName}}
\end{center}

\begin{tcolorbox}[colback=sky, colframe=navy, boxrule=0.9pt, arc=2mm,
    left=2mm, right=2mm, top=1.5mm, bottom=1.5mm]
    \textbf{Primary Objective:} Students will distinguish between chi-square tests for
    homogeneity and independence based on study design, state appropriate hypotheses for
    each test type, and verify the conditions required for two-way table chi-square tests.
    (Big Idea: VAR)
\end{tcolorbox}

\renewcommand{\arraystretch}{1.6}
\begin{skillbox}[Priority Ideas \& Skills]{goldbox}
\begin{minipage}[t]{0.48\linewidth}
    \textbf{AP Skill 1.E --- Identifying an Appropriate Inference Procedure}
    \begin{itemize}
        \item Identify the appropriate chi-square test (homogeneity vs.\ independence)
              based on the study design (VAR-8.J).
        \item Key distinguisher: homogeneity = multiple groups, one categorical variable;
              independence = one sample, two categorical variables.
    \end{itemize}
    \textbf{AP Skill 1.F --- Identifying Hypotheses}
    \begin{itemize}
        \item State $H_0$ and $H_a$ correctly for each test type (VAR-8.I).
        \item Homogeneity: hypothesize about whether distributions are the same across groups.
        \item Independence: hypothesize about whether an association exists.
    \end{itemize}
    \textbf{AP Skill 4.C --- Verifying Conditions}
    \begin{itemize}
        \item Independence condition and large-counts condition for two-way table tests (VAR-8.K).
    \end{itemize}
\end{minipage}
\hfill
\begin{minipage}[t]{0.48\linewidth}
    \textbf{Key Understandings}
    \begin{itemize}
        \item Both tests use the \emph{same} chi-square statistic formula and the
              \emph{same} degrees of freedom: $df = (\text{rows}-1)(\text{cols}-1)$.
              The difference lies entirely in study design, hypothesis wording, and
              conclusion wording.
        \item Test for \textbf{homogeneity}: researcher draws separate random samples from
              two or more populations (or assigns treatments), then records one categorical
              variable. Goal: compare distributions across groups.
        \item Test for \textbf{independence}: researcher draws one random sample from one
              population, then records two categorical variables on each individual.
              Goal: determine whether the two variables are associated.
        \item Conditions: (1) random sampling or random assignment; (2) if sampling without
              replacement, $n \le 10\%N$; (3) all expected counts $> 5$.
    \end{itemize}
\end{minipage}
\end{skillbox}

\begin{skillbox}[{Vocabulary, Concepts, \& Theorems}]{greenbox}
\begin{minipage}[t]{0.48\linewidth}
    \begin{itemize}
        \item \textbf{Test for homogeneity} --- a chi-square test that compares the
              distribution of one categorical variable across two or more populations
              or treatment groups
        \item \textbf{Test for independence} --- a chi-square test that determines
              whether two categorical variables are associated in a single population
        \item \textbf{Association} --- a relationship between two variables such that
              knowing one variable's value changes the probability distribution of the other
    \end{itemize}
\end{minipage}
\hfill
\begin{minipage}[t]{0.48\linewidth}
    \begin{itemize}
        \item \textbf{Stratified random sample} --- a sampling method in which the
              population is divided into subgroups (strata) and a random sample is
              drawn from each stratum; used in homogeneity settings
        \item \textbf{Expected count (two-way table)} ---
              $E = \dfrac{(\text{row total})(\text{column total})}{\text{table total}}$
        \item \textbf{Degrees of freedom (two-way)} ---
              $df = (\text{rows} - 1)(\text{cols} - 1)$
    \end{itemize}
\end{minipage}
\end{skillbox}

\begin{skillbox}[Activate Prior Knowledge \& Spiral Review]{sky}
\begin{minipage}[t]{0.55\linewidth}
    \textbf{Prior Knowledge Check (3--4 min)}
    \begin{itemize}
        \item Lesson 8.4: calculating expected counts from a two-way table using
              $E = (\text{row total} \times \text{column total}) \div n$.
        \item Lesson 8.4: large-counts condition for two-way tables --- all expected
              counts must exceed 5.
        \item Units 6--7: setting up $H_0$ and $H_a$ for significance tests; the
              language of ``no difference'' and ``no association.''
        \item Connect: ``In Lesson 8.4 we computed expected counts and checked
              conditions. Today we figure out \emph{which} test to run and how to
              write the hypotheses correctly.''
    \end{itemize}
    \textit{Warm-up spiral review is attached --- expected counts in a two-way table.}
\end{minipage}
\hfill
\begin{minipage}[t]{0.38\linewidth}
    \textbf{Connections Forward}
    \begin{itemize}
        \item Today: set up the test (identify, hypothesize, conditions)
        \item Lesson 8.6: carry out the chi-square test for homogeneity/independence
              (compute statistic, $p$-value, conclusion)
        \item Lesson 8.7: choose the right chi-square test on the AP exam
    \end{itemize}
\end{minipage}
\end{skillbox}

\begin{skillbox}[Hook]{sky}
    \textbf{Hook (4 min)}\\
    Present two newspaper headlines side by side:
    \begin{itemize}
        \item \textit{``Do students at different schools participate in extracurriculars
              at different rates? Researchers surveyed 100 students at each of three schools.''}
        \item \textit{``Is political affiliation linked to policy opinion? A Gallup poll of
              300 adults reveals a surprising pattern.''}
    \end{itemize}
    Ask: \textbf{``Are these questions the same kind of question, or different kinds?
    What is each study really asking?''} Let students argue. Reveal: both call for a
    chi-square test --- but there are two flavors, and picking the wrong one costs
    points on the AP exam.
\end{skillbox}

\begin{skillbox}[Lesson]{sky}
\begin{minipage}[t]{0.48\linewidth}
    \textbf{Two Study Designs (8 min)}
    \begin{enumerate}
        \item \textbf{Multiple groups, one variable.} A researcher samples from 3 schools
              and records each student's extracurricular choice (Sports/Arts/Neither).
              There are three separate samples and \emph{one} categorical variable.
              $\Rightarrow$ \textbf{Test for Homogeneity.}
        \item \textbf{One group, two variables.} A researcher takes one sample of 300 adults
              and records \emph{two} things about each: political affiliation AND policy view.
              $\Rightarrow$ \textbf{Test for Independence.}
    \end{enumerate}
    \textbf{Decision questions to ask:}
    \begin{itemize}
        \item Were data collected from \emph{one} sample or \emph{multiple} groups?
        \item Were \emph{one} or \emph{two} categorical variables measured per individual?
    \end{itemize}
\end{minipage}
\hfill
\begin{minipage}[t]{0.48\linewidth}
    \textbf{Side-by-Side Comparison}

    \renewcommand{\arraystretch}{1.8}
    \begin{tabularx}{\linewidth}{@{} l X X @{}}
        & \textbf{Homogeneity} & \textbf{Independence} \\
        \hline
        \textbf{Design} & Multiple groups; one variable & One sample; two variables \\
        \textbf{$H_0$} & Distributions of [var] are the same across all [groups] & No association between [var 1] and [var 2] in the population \\
        \textbf{$H_a$} & Distributions differ (at least one group) & There is an association \\
        \textbf{Conclusion} & Refers to distributions across groups & Refers to association in the population \\
        \textbf{df} & $(r-1)(c-1)$ & $(r-1)(c-1)$ \\
        \textbf{Statistic} & $\chi^2 = \sum\dfrac{(O-E)^2}{E}$ & same \\
    \end{tabularx}
\end{minipage}
\end{skillbox}

\begin{skillbox}[Lesson (cont.)]{sky}
\begin{minipage}[t]{0.48\linewidth}
    \textbf{Hypotheses Wording (5 min)}\\[4pt]
    \textbf{Homogeneity example (School survey):}
    \begin{itemize}
        \item $H_0$: The distribution of extracurricular participation is the same
              across Schools A, B, and C.
        \item $H_a$: The distribution of extracurricular participation is not the
              same across all three schools (at least one school differs).
    \end{itemize}
    \textbf{Independence example (Political poll):}
    \begin{itemize}
        \item $H_0$: There is no association between political affiliation and policy
              views in the adult population.
        \item $H_a$: There is an association between political affiliation and policy
              views in the adult population.
    \end{itemize}
    \textit{AP tip: the null always says ``no difference'' (homogeneity) or
    ``no association'' (independence).}
\end{minipage}
\hfill
\begin{minipage}[t]{0.48\linewidth}
    \textbf{Conditions (5 min)}\\[4pt]
    \textbf{1. Independence condition}
    \begin{itemize}
        \item Homogeneity: stratified random sample or randomized experiment.
        \item Independence: simple random sample from the population.
        \item If sampling without replacement: $n \le 10\%N$.
    \end{itemize}
    \textbf{2. Large-counts condition}
    \begin{itemize}
        \item All \emph{expected} counts $> 5$.
        \item Expected count formula: $E = \dfrac{(\text{row total})(\text{col total})}{n}$.
        \item Must check \emph{every} cell in the table.
    \end{itemize}
    \textit{Remind students: conditions are checked the same way regardless of which test.}
\end{minipage}
\end{skillbox}

\begin{skillbox}[Active Monitoring]{redbox}
\begin{minipage}[t]{0.48\linewidth}
    \textbf{Circulate and Check}
    \begin{itemize}
        \item Most common error: labeling a test as ``independence'' simply because
              two variables appear in a table. Redirect: ask ``How many random samples
              were drawn?''
        \item Check hypothesis wording --- $H_0$ for homogeneity must name the
              \emph{distribution}, not just ``the proportions.''
        \item Listen for students mixing up ``same distribution'' with
              ``independent.'' These are distinct statistical claims.
    \end{itemize}
\end{minipage}
\hfill
\begin{minipage}[t]{0.48\linewidth}
    \textbf{Cold-Call Prompts}
    \begin{itemize}
        \item ``A researcher randomly selects 50 men and 50 women and asks each
              their favorite genre of music. Is this homogeneity or independence?
              How do you know?''
        \item ``A researcher takes one SRS of 200 teenagers and asks each their
              favorite genre AND how many hours they stream per day. Homogeneity
              or independence?''
        \item ``What is $df$ for a $3 \times 4$ two-way table?'' (Answer: 6)
    \end{itemize}
\end{minipage}
\end{skillbox}

\begin{skillbox}[Group Work \& Differentiation]{redbox}
\begin{minipage}[t]{0.48\linewidth}
    \textbf{Activity (10--12 min)}\\
    Students work through both Scenario A (three schools, extracurricular) and
    Scenario B (300 adults, political affiliation and policy view). For each:
    \begin{enumerate}
        \item Identify the test type and justify from the study design.
        \item Write $H_0$ and $H_a$ in context.
        \item Identify how independence was achieved.
        \item Calculate 2 expected counts and check the large-counts condition.
    \end{enumerate}
\end{minipage}
\hfill
\begin{minipage}[t]{0.48\linewidth}
    \textbf{Tier R --- Remediate}
    \begin{itemize}
        \item Provide the decision questions as a checklist.
        \item Provide the expected-count formula already filled in for one cell; students
              replicate for one more cell.
        \item Focus on distinguishing ``one sample or multiple?'' rather than the full
              condition check.
    \end{itemize}
    \textbf{Tier A --- Approaching Proficiency}
    \begin{itemize}
        \item Complete all four parts for both scenarios independently.
        \item Verify the large-counts condition by computing all expected counts.
    \end{itemize}
    \textbf{Tier E --- Extension}
    \begin{itemize}
        \item Write a third scenario of their own creation for each test type, swap
              with a partner, and check each other's hypotheses and conditions.
        \item Identify what would change (only the conclusion wording) if we ran
              a homogeneity test vs.\ an independence test on the same data table.
    \end{itemize}
\end{minipage}
\end{skillbox}

\begin{skillbox}[Individual Work \& Assessment]{redbox}
\begin{minipage}[t]{0.48\linewidth}
    \textbf{Exit Ticket (5 min)}
    \begin{enumerate}
        \item A school district randomly selects 80 students from each of 4 grade
              levels and asks which subject (Math/Science/English) is their favorite.
              (a) Identify the correct test. (b) Write $H_0$ in words.
        \item A random sample of 200 adults is asked their marital status
              (Married/Single/Other) and whether they exercise regularly (Yes/No).
              (a) Identify the correct test. (b) Write $H_0$ in words.
    \end{enumerate}
\end{minipage}
\hfill
\begin{minipage}[t]{0.48\linewidth}
    \textbf{Assessment Note}\\
    Collect and sort into three piles:
    \begin{itemize}
        \item Correct test + correct $H_0$ wording $\to$ ready for Lesson 8.6.
        \item Correct test, incorrect $H_0$ wording $\to$ reteach hypothesis language
              at the start of 8.6.
        \item Incorrect test $\to$ pull for a 2-min small-group review of the
              ``one sample vs.\ multiple groups'' decision rule before 8.6.
    \end{itemize}
\end{minipage}
\end{skillbox}

\begin{skillbox}[Reinforcement \& Extension]{goldbox}
\begin{minipage}[t]{0.48\linewidth}
    \textbf{Homework Overview}
    \begin{itemize}
        \item Problem 1 (three parts): classify scenarios as homogeneity or independence
              and write hypotheses.
        \item Problem 2: address the common misconception that ``two variables = independence.''
        \item Problem 3: write out full conditions for a homogeneity study.
        \item Problem 4: AP-style MC on identifying the correct chi-square test.
    \end{itemize}
\end{minipage}
\hfill
\begin{minipage}[t]{0.48\linewidth}
    \textbf{Extension / Preview}
    \begin{itemize}
        \item \textit{Extension}: Find a real news article that reports a chi-square test.
              Identify which type of test was used and whether the hypotheses are stated
              correctly.
        \item \textit{Preview}: Lesson 8.6 carries out the full chi-square test on these
              scenarios --- computing the test statistic, finding the $p$-value, and
              writing a complete conclusion in context.
    \end{itemize}
\end{minipage}
\end{skillbox}

\end{document}
